\chapter{מבוא: מהו סטארט-אפ ולמה רובם נכשלים}

\term{startup} אינו סתם עסק קטן. זהו ארגון זמני שנועד לחפש מודל עסקי
\en{(business model)} בר-חזרה והדיר, בתנאי אי-ודאות קיצונית~\cite{blank2013startup}.
ההבדל מהותי: בעל מסעדה מפעיל מודל ידוע, בעוד שיזם סטארט-אפ \emph{מחפש} מודל
שעדיין אינו קיים. מסיבה זו, סטארט-אפ הוא קודם כול \emph{ניסוי} — ולא תוכנית
ביצוע.

הסטטיסטיקה איננה מחמיאה: כשלושה מכל ארבעה סטארט-אפים מגובי הון אינם מחזירים
את ההשקעה, והסיבה השכיחה ביותר לכישלון אינה טכנולוגיה כושלת אלא היעדר ביקוש
אמיתי בשוק~\cite{cbinsights2021}. כלומר, רוב הצוותים בונים מוצר נכון מבחינה
הנדסית — לבעיה הלא-נכונה.

\section{מסע המזומנים של הסטארט-אפ}

הדרך מרעיון לחברה רווחית עוברת כמעט תמיד דרך ``עקומת ה-\en{J}'': תזרים המזומנים
המצטבר צולל תחילה אל מתחת לאפס (שלב השריפה, \term{burn}), מגיע לשפל, ורק לאחר
שההכנסות עוקפות את העלויות הקבועות הוא מטפס חזרה אל הרווחיות. איור~\ref{fig:jcurve}
ממחיש זאת.

\begin{figure}[h]
  \centering
  \includegraphics[width=0.78\textwidth]{jcurve.pdf}
  \caption{עקומת ה-\en{J} של תזרים המזומנים המצטבר לאורך 24 חודשים. נקודת השפל
  מסמנת את הרגע שבו דרוש הון נוסף כדי לשרוד עד להתאוששות.}
  \label{fig:jcurve}
\end{figure}

\section{תזת הספר}

ספר זה טוען טענה אחת מרכזית: בניית סטארט-אפ נכון היא \emph{תהליך למידה ממושמע},
לא הימור חד-פעמי. נציג שיטה מסודרת — מאימות הבעיה, דרך לולאת
\term{Build--Measure--Learn}, התאמת מוצר-שוק, בניית הצוות וכלכלת היחידה, ועד
ניהול הסיכונים והטעויות הנפוצות. כל פרק מסתמך על עקרונות שהוכחו בספרות
המקצועית~\cite{ries2011lean,thiel2014zero}.
